% =====================================================================
%  4. ÉTUDE DE CAS — RECODAGE DE ROOMEQ
% =====================================================================
\section{Étude de cas : recodage de RoomEq}\label{sec:etude}

\subsection{Contexte et protocole expérimental}\label{sec:protocole-etude}

RoomEq est un égaliseur acoustique temps réel pour Windows et macOS~: il
mesure la réponse d’une pièce par balayage logarithmique, calcule une
correction, l’applique en temps réel sur le son du système, et offre une
console de mixage complète — pistes, égaliseur par piste, routage par carte
son virtuelle, pont ASIO vers les stations de travail audio. La base de
départ est la version 2.2.2, \emph{scellée} au sens du protocole~: balise
git \code{deff25c}, 82 fichiers C\#, 14\,891 lignes de code, 228 tests
unitaires, trois ans de développement intermittent par le développeur et
ses prédécesseurs d’agents. Le candidat cumule les attributs d’un cas
exigeant~: temps réel audio, interactions noyau (pilote de carte virtuelle,
installateur), multi-plateforme, et un patrimoine de défauts d’installation
documentés.

L’étude poursuit trois questions de recherche. \textbf{RQ1}~: la discipline
des trois couches est-elle applicable de bout en bout à une application
réelle complète, sans exception ni accommodement~? \textbf{RQ2}~: quel est
l’effet mesurable de cette discipline sur la défense (tests, avertissements)
et la structure (taille des unités, testabilité) du code~?
\textbf{RQ3}~: que révèle la vérification de l’oracle que le
développement classique ne révèle pas~? Le design est une \emph{étude de
cas unique} au sens de Yin \cite{yin2018}, conduite et rapportée selon les
recommandations de Runeson et H{\"o}st \cite{runeson2009}~; l’unité
d’analyse est le recodage complet, de la base scellée au cristal final. Les
limites inhérentes à ce design — absence de condition témoin, validité
externe — sont traitées en section~\ref{sec:menaces} plutôt que tues.

L’instrumentation est triple. Premièrement, une intégration continue
exécutant construction et tests sur trois systèmes d’exploitation (Linux,
Windows, macOS), activée à partir de la cristallisation finale.
Deuxièmement, des empreintes cryptographiques des données embarquées
(bibliothèques de profils), reproductibles par quiconque exécute la
commande de hachage sur les deux dépôts. Troisièmement, un audit de parité
conduisant famille par famille — les quatorze familles de fonctionnalités
du fichier de lecture de la base scellée — avec, pour chacune, les modules
correspondants, la méthode de vérification et le verdict. Aucune de ces
mesures n’a été conçue après coup~: elles procèdent de la règle R5 elle-même,
appliquée à chaque vague.

\subsection{Application de la méthode}\label{sec:application}

Le recodage s’est déroulé en sept vagues, chacune close par une sauvegarde
complète (dépôt git local comme source de vérité, miroir chiffré Drive avec
bundle d’historique restaurable, miroir GitHub privé) et par la vérification
complète~: construction à 0 erreur et 0 avertissement, 100~\% des tests
verts. L’instanciation des couches est directe~: la couche $\IB$ est
constituée de 44 spécifications en texte structuré, une par module, écrites
\emph{avant} le code~; la couche $\BB$ est le code source versionné, scellé
par balises~; la couche $\CB$ est la carte du dépôt — un index à adresses
stables \code{M-<DOM>-NN}, un statut par module~; le journal $\DB$ du dépôt
compte quinze décisions, chacune avec ses alternatives rejetées. L’oracle
est la suite de tests \emph{portée puis étendue}~: la suite 2.2.2 encode le
comportement réellement validé du terrain, et chaque constante numérique de
référence (tolérances, amplitudes, fréquences) en a été extraite — la
décision D5 du dépôt pose explicitement que les tests sont le contrat
exécutable de la couche I.

Le tableau~\ref{tab:regles} (section~\ref{sec:regles}) a servi de table de
correspondance normative tout au long des vagues~: chaque règle du
protocole y trouvait son homologue concret dans le dépôt. Deux incidents
méritent d’être consignés car ils éprouvent la gouvernance plutôt que le
code. D’une part, l’environnement d’exécution a été \emph{réinitialisé} en
cours de projet — répertoires de travail et journal effacés~; la
restauration complète depuis les sauvegardes n’a coûté aucune perte de
travail, validant en conditions réelles la redondance exigée par la
gouvernance. D’autre part, une demande du développeur en fin de
démonstration — fusionner les modules et réécrire « en natif machine » —
a été refusée par décision motivée D15 après mesure~: l’analyse a établi
que la fusion détruirait les adresses stables et les spécifications par
responsabilité sans aucun gain d’exécution, le fil audio tourne déjà à
152~$\times$ le temps réel avec zéro allocation par bloc
(section~\ref{sec:resultats}). Ce refus est l’illustration la plus directe
de la fonction objectif lexicographique (définition~\ref{def:lex})~: la
demande promettait un gain d’\emph{optimisation} au prix d’une perte de
\emph{sécurité} de gouvernance — inadmissible quel que soit le gain.

\subsection{Observations qualitatives~: trois comparatifs
avant/après}\label{sec:comparatifs}

\subsubsection{Le seuil magique}\label{sec:seuil}

L’analyseur de balayage de la base scellée rejetait les captures invalides
au-dessus d’un seuil pic/médiane de~8 — un littéral sans nom, dupliqué dans
le message d’erreur, sans justification ni test. Le défaut était réel~: la
déconvolution divise par le spectre du balayage au carré, ce qui agit
comme un différenciateur hors bande~; une capture de \emph{bruit seul}
produit un faux pic à rapport mesuré d’environ 186 — largement au-dessus
de~8. Concrètement, un utilisateur ayant mal routé sa capture se voyait
attribuer une « correction » aléatoire au lieu du message d’aide. Le
protocole ne s’est pas contenté de corriger~: il a d’abord établi les faits
par diagnostic empirique (bruit blanc pur~: rapport 186~; bruit passe-bas
réaliste~: 9,2~; balayage réel même atténué~: au-delà de $10^{5}$, mesuré à
$1{,}2 \times 10^{7}$), puis consigné la décision D8 — un seuil de 1000
rejette le bruit seul avec quatre décades de marge et conserve quatre
décades de marge aux balayages faibles.

Le listing~\ref{lst:seuil} montre la conversion. Le défaut est désormais
verrouillé à quatre niveaux~: constante nommée dans le code, justification
documentée, contrat dans la spécification I du module, test exécutable qui
rejette le bruit seul.

\begin{lstlisting}[float=t, caption={Le seuil de détection, avant (haut) et
après (bas). Extraits abrégés des deux dépôts.}, label={lst:seuil}]
// AVANT — v2.2.2 : seuil littéral, sans nom, sans test :
if (peakMag < Math.Max(median * 8.0, 1e-9)) {
    throw new InvalidOperationException(
        "Balayage non détecté ... rapport ... < 8 ...");
}
// APRÈS — v3 : constante nommée, rationale, message dérivé :
/// Seuil pic/médiane. v2.2.2 : 8 laissait passer le bruit
/// seul (ratio ~186, effet différenciateur du 1/|X|^2) ;
/// balayage réel mesuré ~1,2e7 -> seuil 1000 : quatre
/// décades de marge de chaque côté.
public const double DetectionSnr = 1000.0;
if (peakMag < Math.Max(median * DetectionSnr, 1e-9)) {
    ... // message dérivé de la constante
}
\end{lstlisting}

\subsubsection{Le monolithe}\label{sec:monolithe}

Le modèle-vue de la table de mixage concentrait cinq responsabilités —
routage de la carte virtuelle, synchronisation des volumes, installation du
pilote, persistance, politique de préréglages — entrelacées sur 1\,075
lignes~; 39 propriétés observables et commandes y côtoyaient les
dictionnaires de routage et les minuteries. Aucun test ne peut s’approcher
d’une telle classe sans périphérique audio réel. La spécification I de
l’onglet mixage exigeait la décomposition \emph{avant} que le code
n’existe~; le recodage l’a réalisée en cinq fichiers à responsabilité
unique (399~+ 224~+ 209~+ 160~+ 138 lignes) et extrait la politique de
préréglages vers une classe dédiée derrière deux interfaces étroites.

Le gain n’est pas cosmétique. La décomposition a rendus visibles — et
donc corrigeables — deux défauts fonctionnels que le monolithe masquait~:
le message d’aide ASIO n’était jamais affiché sans pilote (la propriété
concernée restait nulle, aucun changement n’était notifié)~; le garde-fou
de liaison d’un préréglage à une sortie laissait passer le préréglage plat
par défaut, déclenchant une demande de nom puis un refus en cascade. Les
deux correctifs sont couverts par des tests, et la politique extraite est
vérifiée \emph{sans périphérique} — ce qui était impensable sur le
monolithe est devenu une suite de tests ordinaire. C’est la règle R3 — le
balisage — appliquée à l’architecture~: un module, une responsabilité, une
adresse stable dans la carte.

\subsubsection{Le caractère invisible}\label{sec:caractere}

Les magasins de réglages sécurisent les noms de fichiers avant d’écrire sur
disque. La base scellée déléguait cette sécurité au système hôte~: la liste
renvoyée par \code{Path.GetInvalidFileNameChars()} est celle de l’OS qui
exécute le code. Le défaut est invisible sur la machine de
développement Windows — et c’est précisément le problème~: un nom sécurisé
sous Linux, où les deux-points sont légaux, peut produire un fichier
invalide s’il est copié sous Windows. Il n’a été révélé que parce que la
vérification R5 de la démonstration exécute les tests sous Linux~:
l’oracle a fait son travail. Le correctif fige l’union des caractères
interdits sur au moins une plateforme cible, indépendamment de l’hôte~; le
contrat « un nom sécurisé est portable » est désormais écrit dans
la spécification I et dans le code. Trois leçons générales s’en
dégagent~: la vérification doit s’exécuter dans l’environnement qui révèle
le défaut~; un défaut de portabilité est un défaut de contrat~; et la
taille du correctif (six lignes remplacées par quinze) est sans rapport
avec son changement de classe — d’un comportement dépendant de l’hôte à
une garantie multi-plateformes par construction.

\subsection{Résultats quantitatifs}\label{sec:resultats}

Le tableau~\ref{tab:mesures} confronte les deux révisions sur les mesures
structurelles, défensives et de gouvernance. Le volume de code est
volontairement resté comparable — le recodage n’a pas gonflé
l’application, il a déplacé l’effort vers la spécification et la
vérification.

\begin{table}[t]
\centering\small
\caption{Confrontation des deux révisions.}
\label{tab:mesures}
\begin{tabularx}{\columnwidth}{@{}Y l l@{}}
\toprule
\textbf{Mesure} & \textbf{v2.2.2} & \textbf{v3} \\
\midrule
Fichiers C\# (\code{src/}) & 82 & 93 \\
Lignes de code (\code{src/}) & 14\,891 & 15\,438 \\
Lignes de tests & 4\,032 & 5\,963 \\
Tests unitaires & 228 & \textbf{437} \\
Plus gros fichier & 1\,075 l. & \textbf{497} l. \\
Plus gros modèle-vue & 1\,075 l. & \textbf{399} l. \\
Avertissements construction & tolérés & \textbf{0 / 0} \\
Spécifications d’intention & 0 & \textbf{44} \\
Modules cartographiés & 0 & \textbf{42} \\
Décisions journalisées & aucune & \textbf{D1–D15} \\
\bottomrule
\end{tabularx}
\end{table}

La suite de tests mérite le détail, car elle est le cœur de la
vérification. Les 228 tests de la base ont été portés intégralement —
c’était l’oracle — puis étendus~: la couche cœur passe de 150 à 220 tests
(parité des données embarquées, comportements aux limites, diagnostics
injectables), la couche audio de 78 à 142 (fournisseurs, bus principal,
anti-larsen vérifiable sous Linux, rognage exact), et une suite
applicative de 75 tests apparaît, rendue possible par la décomposition des
modèles-vues. La non-régression se démontre famille par famille~: les
quatorze familles de fonctionnalités du fichier de lecture de la base
scellée ont été auditées point par point — parité intégrale, 14/14. Les
données embarquées sont contrôlées par empreintes~: la bibliothèque de
corrections casques (321 modèles, 58 marques) et les profils de mesure
micro (36 entrées) sont identiques octet pour octet à la base scellée,
avec des sommes md5 reproductibles sur les deux dépôts. Trois écarts
volontaires subsistent — le seuil D8 (défaut connu corrigé par décision
journalisée), des correctifs de défauts de la base chacun couvert par un
test, et la décomposition structurelle sans changement de comportement
observable — et leur existence même est une propriété du protocole~:
aucun n’est silencieux.

La séquence de \emph{publication} du cristal final éprouve le mécanisme
d’échec honnête du protocole. La balise v3.0.0 a échoué en intégration
continue~: trois tests relisaient un journal vivant avec un verrouillage de
fichier incompatibles avec l’écriture sous Windows — défaut invisible sous
Linux et macOS, hérité d’une connaissance de portage perdue entre les
versions. Le correctif v3.0.1, limité aux tests, a rendu la vérification
verte sur les trois systèmes mais échoué à l’empaquetage (un outil absent
des exécutants de la CI, et une archive source trop inclusive). La version
v3.0.2 a porté l’outillage d’empaquetage sur un socle portable et ajouté
un reproducteur permanent exécuté \emph{avant} toute publication. Trois
balises, aucune déplacée~: chaque échec a été corrigé en patch à contrat
stable, conformément à la gouvernance — premier exemple complet du
mécanisme « échec de CI $\to$ patch honnête » que la règle R5 rend
obligatoire.

Enfin, la performance a été mesurée après une demande d’optimisation
« maximale » (décision D15 du dépôt)~: sur une machine à 2~vCPU partagés,
à 48~kHz stéréo par blocs de 10~ms, la chaîne d’égalisation à 30 bandes et
2 canaux traite un bloc en 66~µs — 152~$\times$ le temps réel — avec
\emph{zéro} octet alloué sur 100\,000 blocs consécutifs (1\,000~s de son)~;
le tampon circulaire producteur-consommateur coûte 0,18~µs par bloc~; la
régénération d’un balayage de 12~s prend 10,7~ms. Il n’existait aucun
problème de performance à corriger — conclusion qui fonde le refus motivé
de la réécriture « native machine » évoqué en
section~\ref{sec:application}.

\subsection{Reverse-specs de l’installateur scellé}\label{sec:reverse}

L’histoire du produit est marquée par trois défauts d’installateur
consécutifs (versions 2.2.0 à 2.2.2), dont l’analyse des journaux et des
commits révèle la classe commune~: des contrats non écrits, dupliqués en
plusieurs exemplaires. Plutôt que de recoder l’installateur (ce qui
ajouterait un chantier non validé à un chantier en attente de validation),
la démonstration a écrit ses \emph{spécifications rétro-actives}
(\emph{reverse-specs})~: huit modules M-INS rédigés en relisant ligne à
ligne le code source scellé. Le binaire reste la référence
exploitable~; le contrat, lui, devient diffable, versionné et testable. La
règle R2 — l’intention d’abord — s’applique à l’envers du temps~: le
contrat précède même la décision de recoder.

\begin{table}[t]
\centering\small
\caption{Les trois défauts historiques de l’installateur et les verrous
posés par les reverse-specs.}
\label{tab:reverse}
\begin{tabularx}{\columnwidth}{@{}Y Y@{}}
\toprule
\textbf{Défaut historique} & \textbf{Verrou posé} \\
\midrule
Structure Win32 déclarée de mémoire — 424 octets au lieu de 32 sur x64,
échec \code{0x6F8} (2 copies) & invariant de taille testé, contrat complet
SetupAPI / newdev / cfgmgr32 \\
Signature INF erronée (\code{\$Windows NT\$} au lieu de
\code{\$CHICAGO\$}), échec \code{0xE0000205} & validateur bloquant à la
construction et à l’installation \\
Analyse de la sortie pnputil sur libellés localisés français (3 copies) &
fragments indépendants des libellés, décodage CP850, repli par balayage,
fixtures partagées \\
\bottomrule
\end{tabularx}
\end{table}

Formellement, la reverse-spec est une \emph{extraction de spécification
sonore}~: étant donné un artefact scellé $\BB_{\mathrm{scellé}}$, on
construit $\IB'$ telle que le comportement observable de l’artefact
satisfait $\IB'$ ($\mathrm{Obs}(\BB_{\mathrm{scellé}}) \models \IB'$),
sans exiger la couverture exhaustive des comportements non observés. La
technique étend la couche I au patrimoine existant — situations que la
spécification du protocole n’avait pas prévues — et fige la connaissance~
: les spécifications dérivent du code relu, pas d’une mémoire de session.
Si la vague de recodage de l’installateur est un jour lancée, elle partira
de contrats déjà testés~; si elle ne l’est jamais, le comportement validé
reste spécifié pour quiconque viendra après.
